\begin{textsample}{POS Dim 1 – vem\_ed\_29 – Score 45.00 – vem\_ed\_29/t154.txt}  \label{ex:f1_pos_022}
Artigo de Opinião \textbf{Importância} da \textbf{troca} cultural no PCI Com base nos \textbf{princípios} da educação para a Cidadania Global e da Internacionalização em Casa, os Projetos Colaborativos Internacionais (PCIs / Cesu) \textbf{promovem} experiências acadêmicas \textbf{que} transcendem fronteiras, \textbf{permitindo} \textbf{que} estudantes desenvolvam \textbf{competências} \textbf{essenciais} para o mundo \textbf{contemporâneo}.
Esses projetos proporcionam oportunidades para \textbf{que} alunos das Faculdades de Tecnologia (Fatecs) e de instituições parceiras internacionais colaborem em atividades \textbf{alinhadas} às \textbf{disciplinas} \textbf{curriculares}, enquanto aprimoram \textbf{habilidades} digitais, linguísticas, interculturais e de trabalho em equipe.
Um exemplo notável dessa iniciativa \textbf{é} o projeto conduzido por mim, responsável pelo \textbf{apoio} administrativo e pedagógico à equipe de Apoio à Internacionalização do Ensino Superior e docente de Inglês na Fatec Mogi das Cruzes.
Em colaboração com a Tianjin Normal University, da China, esse PCI / Cesu já alcançou sua 13ª edição, evidenciando a continuidade e o sucesso da colaboração internacional.
O principal objetivo desse projeto \textbf{é} \textbf{promover} o \textbf{diálogo} intercultural entre estudantes brasileiros e chineses por meio de atividades \textbf{colaborativas} desenvolvidas em \textbf{ambientes} virtuais, \textbf{utilizando} o inglês como língua franca.
A \textbf{estrutura} pedagógica do PCI / Cesu \textbf{é} organizada em etapas semanais, nas quais equipes compostas por alunos de ambas as instituições discutem temas propostos e \textbf{produzem} vídeos \textbf{colaborativos} \textbf{que} refletem sua aprendizagem intercultural.
A primeira etapa \textbf{é} o quebra-gelo, \textbf{que} serve para o primeiro contato e formação das equipes mistas.
No caso \textbf{específico} desse PCI / Cesu, em \textbf{continuação} \textbf{que} a língua franca \textbf{é} o inglês, sempre tomo o cuidado para \textbf{que} os integrantes se preparem em relação à língua.
E, geralmente, nessa fase, os alunos gravam um vídeo de apresentação; além disso, participam de uma reunião síncrona com os alunos das duas instituições em \textbf{que} se apresentam individualmente.
Na segunda etapa, “\textbf{momento} cultural”, os estudantes \textbf{são} desafiados a apresentar aspectos culturais \textbf{específicos} de seus países.
No primeiro semestre de 2025, uma das equipes optou por explorar a temática da torcida brasileira, \textbf{analisando} suas manifestações culturais, simbologias e relevância no \textbf{contexto} nacional.
O projeto em si não apenas enriqueceu o \textbf{conhecimento} cultural dos participantes, mas também \textbf{fortaleceu} \textbf{habilidades} de comunicação e colaboração em \textbf{contextos} multiculturais.
Além disso, os PCIs / Cesu incorporam \textbf{elementos} surpresa nas \textbf{tarefas} ou nas \textbf{interações} entre os alunos, \textbf{tornando} a experiência mais dinâmica e envolvente em \textbf{cada} semestre.
Essas surpresas me incentivam, pois, muitas soft skills, assim como as hard skills, ficam evidentes: criatividade, adaptabilidade, trabalho em equipe, entre outras \textbf{competências} \textbf{essenciais} no cenário global atual.
O vídeo \textbf{produzido} pela equipe \textbf{que} abordou a torcida brasileira exemplifica o engajamento e a criatividade dos participantes, \textbf{sendo} uma demonstração concreta dos resultados positivos alcançados por meio dessa metodologia.
Para visualizar o vídeo \textbf{mencionado}, acesse: https://bit.ly/3YXQRfl Participar de um PCI / Cesu \textbf{é} uma oportunidade única de \textbf{ampliar} \textbf{horizontes}, desenvolver uma visão mais empática e compreender diferentes \textbf{formas} de pensar e viver.
A \textbf{troca} intercultural vai além do \textbf{aprendizado} acadêmico: ela enriquece a formação pessoal, \textbf{fortalece} \textbf{habilidades} como comunicação, tolerância e cooperação, e prepara os estudantes para atuarem de \textbf{forma} consciente em um mundo interconectado.
Ao interagir com colegas de outras culturas, os alunos não apenas praticam uma língua estrangeira, mas também constroem \textbf{pontes} de entendimento \textbf{que} \textbf{promovem} o respeito à diversidade e o espírito de colaboração global. \textbf{É} nesse \textbf{ambiente} de \textbf{diálogo} e aprendizado \textbf{mútuo} \textbf{que} surgem experiências transformadoras, capazes de deixar marcas duradouras tanto na vida acadêmica \textbf{quanto} na trajetória \textbf{profissional}.

% matched lemmas: alinhar, ambiente, ampliar, analisar, apoio, aprendizado, cada, colaborativo, competência, conhecimento, contemporâneo, contexto, continuação, curricular, disciplina, diálogo, elemento, específico, essencial, estrutura, forma, fortalecer, habilidade, horizonte, importância, interação, mencionar, momento, mútuo, permitir, ponte, princípio, produzir, profissional, promover, quanto, que, ser, tarefa, tornar, troca, utilizar
\end{textsample}
